\chapter{Connect the Store to PostgreSQL}

The store needs its own schema lifecycle even though it shares the PostgreSQL
jail with the catalog. We will reuse the experiment's existing \texttt{app}
role and password, but create separate development and test databases. This
prevents either Rails application from migrating or clearing the other's data.

\section{Exercise 12A: Create the store databases}

Open the database jail console and create two databases owned by the existing
\texttt{app} role:

\begin{lstlisting}[language=sh]
su -l postgres -c \
  'createdb --owner=app store_development'
su -l postgres -c \
  'createdb --owner=app store_test'
\end{lstlisting}

The first database is for interactive development. The second is disposable:
Rails may create, migrate, and clear it while running tests.

\section{Exercise 12B: Protect the shared password}

Return to the store jail's root console. Create a password file readable by the
store service account:

\begin{lstlisting}[language=sh]
install -o root -g store -m 0640 /dev/null \
  /usr/local/etc/store-db-password
vi /usr/local/etc/store-db-password
\end{lstlisting}

Paste the same experiment-only \texttt{app} password used by the catalog, with
no surrounding quotes. Save the file and verify its ownership and mode:

\begin{lstlisting}[language=sh]
stat -f '%Su:%Sg %Sp %N' \
  /usr/local/etc/store-db-password
\end{lstlisting}

The result should begin with \texttt{root:store -rw-r-----}.

\section{Exercise 12C: Configure both Rails environments}

Edit the generated database configuration as the application user:

\begin{lstlisting}[language=sh]
su -l store
cd /usr/local/www/store
vi config/database.yml
\end{lstlisting}

Replace the generated development and test blocks with:

\begin{lstlisting}
development: &dev
  <<: *default
  host: 10.17.89.23
  port: 5432
  database: store_development
  username: app
  password: <%= File.read(
    "/usr/local/etc/store-db-password").strip %>

test:
  <<: *dev
  database: store_test
\end{lstlisting}

The YAML anchor shares the server, role, and protected password between the
two environments while the explicit test override preserves database
isolation.

\section{Exercise 12D: Verify development and test connections}

Ask Active Record to report the development connection across the jail
boundary:

\begin{lstlisting}[language=sh]
bin/rails runner 'p ActiveRecord::Base.connection.\
select_one("SELECT current_database(), current_user, \
inet_client_addr(), inet_server_addr()")'
\end{lstlisting}

The result should identify:

\begin{itemize}
  \item database \texttt{store\_development};
  \item role \texttt{app};
  \item client \texttt{10.17.89.25}; and
  \item server \texttt{10.17.89.23}.
\end{itemize}

Prepare the isolated test database and run the generated test suite:

\begin{lstlisting}[language=sh]
RAILS_ENV=test bin/rails db:prepare
bin/rails test
\end{lstlisting}

The minimal application may report zero tests at this stage; success means
both commands finish without a database connection or configuration error.

\section{Checkpoint}

Stop after both Rails environments connect successfully. The store now has a
working database boundary but no product page. Next we will make its first
application feature call the catalog API rather than reading catalog data
directly from PostgreSQL.
